% !PRINTED 490
\weberpage{490}

Durch Multiplikation mit $u^{12}v^{12}$ geht die Gleichung (3) in eine Form über, welche $u$, $v$ nicht im Nenner enthält, die wir für den Augenblick mit

\weq{4}{F(u, v) = 0}

bezeichnen, so daß für ein unbestimmtes $x$ und $\omega$:

\weq{5}{F'[f(\omega), x] = [x - f(23\,\omega)] \prod \left[x - f\left(\frac{c+\omega}{23}\right)\right].}
% [sic] printed prime on $F'$ is inconsistent with $F(u,v)$ in (4) and $F(u,u)$ in (6)

Wir untersuchen nun, für welche Werte von $\omega$ die Funktion:

\weq{6}{\begin{aligned}
F(u, u) &= F[f(\omega), f(\omega)] \\
&= [f(\omega) - f(23\,\omega)] \prod \left[f(\omega) - f\left(\frac{c+\omega}{23}\right)\right]
\end{aligned}}

verschwindet. Es verschwindet zunächst offenbar der Faktor:

\[
f(\omega) - f\left(\frac{\omega}{23}\right)
\]
für
\[
u = f(\sqrt{-23}) = f\left(\frac{\sqrt{-23}}{23}\right).
\]

Verstehen wir ferner unter $r$ eine der Zahlen $\pm 2$ oder $\pm 4$, und dementsprechend unter

\[
m = 23 - r^2
\]
entweder 19 oder 7, so verschwinden von den Faktoren von (6) für $\omega = \sqrt{-m}$ je \gesperrt{zwei}, nämlich:

\[
f(\omega) - f\left(\frac{24\,r + \omega}{23}\right),
\]
denn es ist

\[
\begin{aligned}
f\left(\frac{24\,r + \sqrt{-m}}{23}\right) &= e^{\frac{r\pi i}{24}} f\left(\frac{r + \sqrt{-m}}{23}\right) \\
&= e^{\frac{r\pi i}{24}} f\left(\frac{1}{r - \sqrt{-m}}\right) = f(\sqrt{-m}).
\end{aligned}
\]

Wir fügen noch hinzu, daß

\[
\frac{f(\omega) - f\left(\dfrac{24\,r + \omega}{23}\right)}{f(\omega) - f(\sqrt{-m})}
\]
für $\omega = \sqrt{-m}$ einer endlichen Grenze zustrebt, wie man durch Differentiation nach $\omega$ erkennt (\S\,54), und daraus folgt, daß $F(u, u)$ durch $[u - f(\sqrt{-m})]^2$ teilbar ist.
